\chapter{立体几何}
\setcounter{choicecounter}{0}

%% ---- 知识框架 ---- %%
\begin{knowledgeframe}
\begin{itemize}
	\item \textbf{空间线面关系}：线线平行/垂直、线面平行/垂直、面面平行/垂直的判定与性质。
	\item \textbf{充分必要条件}：在立体几何中，线面平行/垂直的条件判定是充要条件分析的重点。
	\item \textbf{面面交线}：$\alpha\cap\beta=m$ 时，直线与两平面的关系可通过交线传递。
	\item \textbf{线面成角}：直线与平面所成的角等于该直线与它在平面内射影所成的角，范围 $[0,\frac{\pi}{2}]$。
\end{itemize}
\end{knowledgeframe}

%% ---- 第一节：空间线面关系 ---- %%
\section{空间线面关系}

\startexercise

\subsection*{A组}

\begin{choice}{2025合肥高三二检}{线面共面与充要条件}
	若空间中三条不同的直线 $a, b, c$ 满足 $a \perp c$, $b \perp c$, 则 $a // b$ 是 $a, b, c$ 共面的
	\begin{tasks}(2)
		\task 充分不必要条件
		\task 必要不充分条件
		\task 充要条件
		\task 既不充分也不必要条件
	\end{tasks}
\end{choice}

\begin{choice}{2025年雅礼月考八}{面面平行的判定}
	设 $\alpha, \beta, \gamma$ 表示三个不同的平面，$m$ 表示直线，则下列选项中，使得 $\alpha // \beta$ 的是
	\begin{tasks}(2)
		\task $m // \alpha, m // \beta$
		\task $m \perp \alpha, m \perp \beta$
		\task $\gamma // \alpha, \gamma // \beta$
		\task $\gamma \perp \alpha, \gamma \perp \beta$
	\end{tasks}
\end{choice}

\begin{choice}{2024年天津卷选择题\#6}{线面平行与垂直}
	若 $m,n$ 为两条不同的直线，$\alpha$ 为一个平面，则下列结论中正确的是
	\begin{tasks}(2)
		\task 若 $m//\alpha$，$n//\alpha$，则 $m\perp n$
		\task 若 $m//\alpha$，$n//\alpha$，则 $m//n$
		\task 若 $m//\alpha$，$n\perp\alpha$，则 $m\perp n$
		\task 若 $m//\alpha$，$n\perp\alpha$，则 $m,n$ 相交
	\end{tasks}
\end{choice}

\subsection*{B组}

\begin{choice}{2024年高考全国甲卷理科选择题\#10}{面面交线与线面关系}
	设 $\alpha$、$\beta$ 是两个平面，$m$、$n$ 是两条直线，且 $\alpha\cap\beta=m$。下列四个命题：\\
	I. 若 $m//n$，则 $n//\alpha$ 或 $n//\beta$。\\
	II. 若 $m\perp n$，则 $n\perp\alpha$ 且 $n\perp\beta$。\\
	III. 若 $n//\alpha$，且 $n//\beta$，则 $m//n$。\\
	IV. 若 $n$ 与 $\alpha$ 和 $\beta$ 所成的角相等，则 $m\perp n$。\\
	其中所有真命题的编号是
	\begin{tasks}(2)
		\task I、III
		\task II、IV
		\task I、II、III
		\task I、III、IV
	\end{tasks}
\end{choice}

\stopexercise
